\documentclass[12pt,a4paper]{article}

\usepackage[spanish]{babel}
\usepackage[utf8]{inputenc}
\usepackage[T1]{fontenc}
\usepackage{geometry}
\usepackage{amsmath}
\usepackage{amssymb}
\usepackage{booktabs}
\usepackage[ruled,vlined,linesnumbered]{algorithm2e}
\usepackage[hidelinks]{hyperref}
\usepackage{float}
\usepackage{graphicx}
\newcommand{\maybeimage}[3]{%
  \IfFileExists{#1}{\includegraphics[#2]{#1}}%
  {\fbox{\parbox{0.7\textwidth}{\centering\vspace{1cm}[Figura pendiente: #3]\vspace{1cm}}}}%
}
\usepackage{caption}
\usepackage{parskip}
\usepackage{setspace}
\usepackage{titlesec}
\usepackage{xcolor}
\usepackage{array}

\geometry{top=2.5cm, bottom=2.5cm, left=2.5cm, right=2.5cm}
\setlength{\parskip}{6pt}

% Ajuste de espaciado en secciones
\titlespacing*{\section}{0pt}{14pt}{6pt}
\titlespacing*{\subsection}{0pt}{10pt}{4pt}

% Comando para valores pendientes
\newcommand{\placeholder}[0]{\textit{---}}

% Configuración de algorithm2e
\SetKwInput{KwIn}{Entrada}
\SetKwInput{KwOut}{Salida}
\SetKw{KwReturn}{devolver}
\SetKwFor{Para}{para}{hacer}{fin para}
\SetKwFor{Mientras}{mientras}{hacer}{fin mientras}
\SetKwIF{Si}{SiNo}{OSiNo}{si}{entonces}{si no si}{si no}{fin si}
\SetKwFunction{fitness}{fitness}
\SetKwFunction{fix}{fix}
\SetKwFunction{createSolution}{crearSolución}
\SetKwFunction{shuffle}{barajar}

%-----------------------------------------------------------------------
\begin{document}

%-----------------------------------------------------------------------
% PORTADA
%-----------------------------------------------------------------------
\begin{titlepage}
  \centering
  \vspace*{2cm}
  {\Large \textbf{Universidad de Granada}\\[0.3cm]
          \textbf{ETSIIT Informática y Telecomunicación}\\[0.3cm]
          Dpto.\ de Ciencias de la Computación e Inteligencia Artificial}

  \vspace{2.5cm}
  {\LARGE \textbf{Metaheurísticas}}\\[0.5cm]
  {\Large Práctica 1}\\[0.3cm]
  {\large Técnicas de Búsqueda Local y Algoritmos Greedy\\
           para el Problema de Portfolio}

  \vspace{1.5cm}
  {\large Curso 2025--2026}

  \vfill
  \begin{tabular}{rl}
    \textbf{Nombre:}  & Jose Angel Carretero Montes \\[4pt]
    \textbf{DNI:}     & 25628892S \\[4pt]
    \textbf{Email:}   & \texttt{joseangel1@correo.ugr.es} \\[4pt]
    \textbf{Grupo:}   & Grupo 3 \\[4pt]
    \textbf{Horario:} & Martes 15:30--17:30 \\
  \end{tabular}
  \vspace{1.5cm}
\end{titlepage}

%-----------------------------------------------------------------------
% ÍNDICE
%-----------------------------------------------------------------------
\tableofcontents
\newpage

%-----------------------------------------------------------------------
\section{Descripción del Problema}
%-----------------------------------------------------------------------

El \textbf{Problema de Portfolio} consiste en determinar la asignación óptima de capital
entre $N$ activos financieros con el objetivo de maximizar el retorno esperado
minimizando simultáneamente el riesgo.

\subsection*{Representación de los datos}

Para cada activo $i \in \{1,\ldots,N\}$ y día $t \in \{1,\ldots,T\}$ se dispone
del retorno diario $\mu_{ti}$, calculado como el ratio de mejora relativa del valor:
\[
  \mu_{ti} = \frac{V_{ti} - V_{(t+1)i}}{V_{ti}}
\]
Los datos empleados cubren el periodo 2015--2024 para entrenamiento y el año 2025
para evaluación de generalización.

\subsection*{Función objetivo}

Dado un vector de pesos $w = (w_1,\ldots,w_N)$, la calidad de la cartera se mide como:
\[
  \text{fitness}(w) = \text{Beneficio}(w) - \lambda \cdot \text{Riesgo}(w)
\]
donde:
\[
  \text{Beneficio}(w) = \sum_{i=1}^{N} w_i \cdot \sum_{t=1}^{T}\log(1+\mu_{ti})
  \qquad
  \text{Riesgo}(w) = \sum_{i=1}^{N}\sum_{j=1}^{N} w_i \cdot \Sigma_{ij} \cdot w_j
\]
siendo $\Sigma_{ij} = \frac{1}{T}\sum_{t=1}^{T}(\mu_{ti}-\bar{\mu}_i)(\mu_{tj}-\bar{\mu}_j)$
la covarianza entre activos $i$ y $j$, y $\lambda = 500$ el factor de penalización
del riesgo.

\subsection*{Restricciones}

\[
  \sum_{i=1}^{N} w_i = 1 \qquad
  w_i = 0 \;\vee\; lo \le w_i \le hi \quad \forall i \in 1\ldots N
\]

Los rangos $[lo, hi]$ varían según el conjunto de datos (véase Sección~\ref{sec:experimentos}).

%-----------------------------------------------------------------------
\section{Aplicación de los Algoritmos al Problema}
%-----------------------------------------------------------------------

\subsection{Representación de soluciones}

Una solución queda representada como un vector $w \in \mathbb{R}^N$
donde cada componente $w_i$ indica la fracción del capital invertida en el activo $i$.
La restricción de cardinalidad implica que $w_i$ puede ser exactamente cero
(no se invierte en esa empresa) o estar en el intervalo $[lo, hi]$.

\subsection{Generación aleatoria de soluciones}

Para construir una solución aleatoria válida se procede como sigue:

\begin{enumerate}
  \item Se elige aleatoriamente un número de empresas en las que \emph{no} se invierte,
        $n_0 \in [1, N{-}1]$.
  \item Los índices restantes $n_+ = N - n_0$ se seleccionan barajando un vector de índices.
  \item A cada empresa activa se le asigna un peso aleatorio en $[lo, hi]$.
  \item Se normaliza: $w_i \leftarrow w_i / \sum_k w_k$.
  \item Se llama al operador \fix{} para garantizar factibilidad.
\end{enumerate}

\subsection{Operador de reparación (\texttt{fix})}

Tras cualquier modificación de la solución puede violarse alguna restricción.
El operador de reparación actúa en tres fases:

\begin{enumerate}
  \item \textbf{Normalización:} $w_i \leftarrow w_i / \sum_k w_k$.
  \item \textbf{Recorte de rango:} si $0 < w_i < lo$ entonces $w_i \leftarrow 0$;
        si $w_i > hi$ entonces $w_i \leftarrow hi$.
  \item \textbf{Redistribución iterativa:} mientras $|\sum w_i - 1| > \varepsilon$:
        \begin{itemize}
          \item Si falta dinero por invertir: se reparte la diferencia entre las empresas
                activas con $w_i < hi$. Si no hay candidatos, se activa una empresa
                inactiva asignándole $lo$.
          \item Si sobra dinero: se retira la diferencia de las empresas con $w_i > lo$.
        \end{itemize}
\end{enumerate}

\subsection{Precómputo de matrices}

Para acelerar la evaluación, se precalculan antes de ejecutar los algoritmos:
el vector de factores de beneficio $b_i = \sum_t \log(1+\mu_{ti})$
y la matriz de covarianza $\Sigma$, tanto para los datos de entrenamiento (2015--2024)
como para los de test (2025). Así, \texttt{fitness} se calcula en $O(N^2)$.

%-----------------------------------------------------------------------
\section{Algoritmos}
%-----------------------------------------------------------------------

\subsection{Algoritmo Aleatorio (\textit{Random Search})}

Genera \texttt{MAX\_EVALS} soluciones aleatorias independientes y devuelve la mejor.
Sirve como referencia inferior de calidad.

\begin{algorithm}[H]
\caption{Búsqueda Aleatoria}
\KwIn{problema, maxEvals}
\KwOut{mejor solución encontrada}
$\text{mejorFit} \leftarrow -\infty$\;
\Para{$i \leftarrow 1$ \KwTo $\text{maxEvals}$}{
  $s \leftarrow$ \createSolution{}\;
  $f \leftarrow$ \fitness{$s$}\;
  \Si{$f > \text{mejorFit}$}{
    $\text{mejorFit} \leftarrow f$\;
    $\text{mejorSol} \leftarrow s$\;
  }
}
\KwReturn{$\text{mejorSol}$}
\end{algorithm}

\subsection{Algoritmo Greedy}

Construye una solución de forma determinista asignando capital en orden decreciente
de una función heurística que estima el valor marginal de cada activo.

\subsubsection*{Función heurística}

La heurística para el activo $i$ estima su contribución neta al fitness si fuese
el único activo de la cartera, penalizando su covarianza media con el resto:
\[
  h(i) = b_i - \lambda \cdot \frac{1}{N}\sum_{j=1}^{N} \Sigma_{ij}
\]
donde $b_i$ es el factor de beneficio precalculado y $\Sigma_{ij}$ la covarianza.

\begin{algorithm}[H]
\caption{Greedy}
\KwIn{problema}
\KwOut{solución construida}
Calcular $h(i)$ para todo $i \in \{1,\ldots,N\}$\;
Ordenar empresas en orden decreciente por $h(i)$ $\rightarrow$ lista \textit{companies}\;
$w \leftarrow \mathbf{0}$;\quad $suma \leftarrow 0$\;
\Mientras{$suma < 1 - \varepsilon$ \textbf{y} quedan empresas en \textit{companies}}{
  $c \leftarrow$ primera empresa de \textit{companies}\;
  Eliminar $c$ de \textit{companies}\;
  $resto \leftarrow 1 - suma$\;
  $w[c] \leftarrow \min(hi,\; resto)$\;
  $suma \leftarrow suma + w[c]$\;
}
\fix{$w$}\tcp*{redistribuye si quedan restos por rounding}
\KwReturn{$w$}
\end{algorithm}

\noindent El greedy es determinista: produce siempre la misma solución independientemente
de la semilla, por lo que su resultado es idéntico en las 50 ejecuciones.

\subsection{Búsqueda Local (BL)}

Implementa una búsqueda de \textbf{primero mejor} (\textit{first improvement})
sobre el vecindario definido por todos los pares de empresas $(p_1, p_2)$,
$p_1 \neq p_2$, con $p_1, p_2 \in \{1,\ldots,N\}$.

\subsubsection*{Operador de generación de vecino}

Dado un par $(origen, destino)$, el vecino se obtiene transfiriendo un
$ratio = 0{,}4$ del peso del activo origen al activo destino:
\[
  \text{change}(w, origen, destino) : \quad
  w'_{origen} = w_{origen} \cdot (1 - ratio),\quad
  w'_{destino} = w_{destino} + ratio \cdot w_{origen}
\]
Tras el cambio se aplica \fix{} para restaurar la factibilidad.

\subsubsection*{Esquema de la BL}

\begin{algorithm}[H]
\caption{Búsqueda Local (primero mejor)}
\KwIn{problema, maxEvals}
\KwOut{mejor solución encontrada}
$P \leftarrow$ todos los pares $(i,j)$ con $i \neq j$, $i,j \in \{1,\ldots,N\}$\;
$s \leftarrow$ \createSolution{};\quad $f \leftarrow$ \fitness{$s$};\quad $evals \leftarrow 1$\;
$mejora \leftarrow \text{verdadero}$\;
\Mientras{$mejora$ \textbf{y} $evals < maxEvals$}{
  $mejora \leftarrow \text{falso}$\;
  \shuffle{$P$}\;
  \Para{cada $(origen, destino) \in P$}{
    \Si{$evals \geq maxEvals$}{\textbf{romper}\;}
    $s' \leftarrow s$\;
    $s'[origen] \leftarrow s'[origen] \cdot 0{,}6$\;
    $s'[destino] \leftarrow s'[destino] + s[origen] \cdot 0{,}4$\;
    \fix{$s'$}\;
    $f' \leftarrow$ \fitness{$s'$}\;
    $evals \mathrel{+}= 1$\;
    \Si{$f' > f$}{
      $s \leftarrow s'$;\quad $f \leftarrow f'$;\quad $mejora \leftarrow \text{verdadero}$\;
      \textbf{romper}\;
    }
  }
}
\KwReturn{$s$}
\end{algorithm}

\noindent El vecindario tiene $N \cdot (N{-}1)$ vecinos. La exploración se baraja en
cada iteración para evitar sesgos de orden. La condición de parada es la convergencia
(ninguna mejora en un recorrido completo) o el agotamiento del presupuesto.


%-----------------------------------------------------------------------
\section{Estructura del Código y Manual de Usuario}
%-----------------------------------------------------------------------

\subsection{Estructura de directorios}

\begin{verbatim}
portfolio/
  main.cpp                  # Punto de entrada, experimentos de 50 semillas
  CMakeLists.txt            # Sistema de construcción (CMake)
  boxplot.py                # Script Python para boxplots (pandas + seaborn)
  LEEME                 # Manual de usuario y descripción del proyecto
  data/
    ibex_35.csv             # 30 empresas IBEX 35, 2015-2025
    syp_100.csv             # 97 empresas S&P 100, 2015-2025
    syp_500.csv             # 457 empresas S&P 500, 2015-2025
  common/
    mh.h                    # Interfaz genérica MH<tDomain>
    problem.h               # Interfaz genérica Problem<tDomain>
    solution.h              # Typedefs tSolution, tFitness
    random.hpp              # Generador pseudoaleatorio de alta calidad
    util.h                  # Utilidades auxiliares
  inc/
    portfolioProblem.h      # Declaración del problema de portfolio
    randomSearch.h          # Algoritmo aleatorio
    greedy.h                # Algoritmo greedy
    localsearch.h           # Búsqueda local (primero mejor)
  src/
    portfolioProblem.cpp    # Carga de datos, precómputo, fitness, fix
    randomSearch.cpp        # Implementación búsqueda aleatoria
    greedy.cpp              # Implementación greedy
    localsearch.cpp         # Implementación búsqueda local
\end{verbatim}

\subsection{Compilación}

Desde el directorio raíz del proyecto:

\begin{verbatim}
cmake .
make
\end{verbatim}

\subsection{Ejecución}

El ejecutable \texttt{main} acepta opcionalmente dos parámetros:

\begin{verbatim}
./main [max_evals] [n_seeds]
\end{verbatim}

\noindent Ejemplos de ejecución:
\begin{verbatim}
./main                  # valores por defecto: 10000 evaluaciones, 50 semillas
./main 5000             # 5000 evaluaciones, 50 semillas
./main 10000 30         # 10000 evaluaciones, 30 semillas
\end{verbatim}

Imprime las tablas de medias por pantalla y genera los ficheros CSV:
\texttt{results\_ibex\_35.csv}, \texttt{results\_syp\_100.csv},
\texttt{results\_syp\_500.csv}.

Para generar los boxplots (requiere Python con \texttt{pandas} y \texttt{seaborn}):

\begin{verbatim}
python3 boxplot.py
\end{verbatim}

Las semillas utilizadas son $1, 2, \ldots, \texttt{n\_seeds}$, iguales
para todos los algoritmos y conjuntos de datos, garantizando comparaciones justas.

%-----------------------------------------------------------------------
\section{Experimentos y Resultados}
\label{sec:experimentos}
%-----------------------------------------------------------------------

\subsection{Casos del problema}

\begin{center}
\begin{tabular}{lrrrl}
\toprule
\textbf{Fichero} & \textbf{Empresas ($N$)} & \textbf{Días train ($T$)} & \textbf{Días test} & \textbf{Rango $[lo, hi]$} \\
\midrule
\texttt{ibex\_35.csv}  & 30  & $\sim$2560 & 254 & $[0.005,\; 0.08]$ \\
\texttt{syp\_100.csv}  & 97  & $\sim$2516 & 249 & $[0.005,\; 0.05]$ \\
\texttt{syp\_500.csv}  & 457 & $\sim$2516 & 249 & $[0.005,\; 0.02]$ \\
\bottomrule
\end{tabular}
\end{center}

Parámetros comunes: $\lambda = 500$, \texttt{MAX\_EVALS} $= 10\,000$,
50 semillas ($1$--$50$).

\subsection{Resultados}

Las tablas muestran la media de las 50 ejecuciones. El fitness reportado para el
Greedy es siempre el mismo al ser determinista.

\begin{table}[H]
\centering
\caption{Resultados IBEX 35 --- $N=30$, $T=2559$, $[lo,hi]=[0.005, 0.08]$ (medias, 50 semillas)}
\begin{tabular}{lrrrrr}
\toprule
\textbf{Algoritmo} & \textbf{Fit.\ 2015--24} & \textbf{Fit.\ 2025} & \textbf{Ben.\ 2025} & \textbf{Evals} & \textbf{Tiempo (s)} \\
\midrule
Greedy  & $-0.2128$ & $0.0907$ & $0.1152$ & $1$         & $0.000$ \\
Random  & $-0.5833$ & $0.0910$ & $0.1180$ & $10\,000$   & $0.020$ \\
BL      & $-0.2128$ & $0.0907$ & $0.1152$ & $1\,852$  & $0.001$ \\
\bottomrule
\end{tabular}
\end{table}

\begin{table}[H]
\centering
\caption{Resultados S\&P 100 --- $N=97$, $T=2516$, $[lo,hi]=[0.005, 0.05]$ (medias, 50 semillas)}
\begin{tabular}{lrrrrr}
\toprule
\textbf{Algoritmo} & \textbf{Fit.\ 2015--24} & \textbf{Fit.\ 2025} & \textbf{Ben.\ 2025} & \textbf{Evals} & \textbf{Tiempo (s)} \\
\midrule
Greedy  & $0.8666$ & $-0.0407$ & $0.0099$ & $1$         & $0.000$ \\
Random  & $0.4342$ & $0.0065$  & $0.0518$ & $10\,000$   & $0.050$ \\
BL      & $0.8665$ & $-0.0393$ & $0.0111$ & $10\,000$   & $0.029$ \\
\bottomrule
\end{tabular}
\end{table}

\begin{table}[H]
\centering
\caption{Resultados S\&P 500 --- $N=457$, $T=2516$, $[lo,hi]=[0.005, 0.02]$ (medias, 50 semillas)}
\begin{tabular}{lrrrrr}
\toprule
\textbf{Algoritmo} & \textbf{Fit.\ 2015--24} & \textbf{Fit.\ 2025} & \textbf{Ben.\ 2025} & \textbf{Evals} & \textbf{Tiempo (s)} \\
\midrule
Greedy  & $1.0430$ & $-0.1157$ & $-0.0628$ & $1$         & $0.000$ \\
Random  & $0.2406$ & $-0.0390$  & $0.0061$  & $10\,000$   & $1.193$ \\
BL      & $0.9491$ & $-0.0986$ & $-0.0484$ & $10\,000$   & $0.745$ \\
\bottomrule
\end{tabular}
\end{table}

\subsection{Boxplots}

Las Figuras~\ref{fig:box_ibex}--\ref{fig:box_syp500} muestran la distribución del
fitness de entrenamiento en las 50 ejecuciones para cada conjunto de datos.

\begin{figure}[H]
  \centering
  \maybeimage{../../results/all_algorithms/results_ibex_35.png}{width=0.7\textwidth}{results\_ibex\_35.png}
  \caption{Distribución del fitness (2015--2024) en IBEX 35.}
  \label{fig:box_ibex}
\end{figure}

\begin{figure}[H]
  \centering
  \maybeimage{../../results/all_algorithms/results_syp_100.png}{width=0.7\textwidth}{results\_syp\_100.png}
  \caption{Distribución del fitness (2015--2024) en S\&P 100.}
  \label{fig:box_syp100}
\end{figure}

\begin{figure}[H]
  \centering
  \maybeimage{../../results/all_algorithms/results_syp_500.png}{width=0.7\textwidth}{results\_syp\_500.png}
  \caption{Distribución del fitness (2015--2024) en S\&P 500.}
  \label{fig:box_syp500}
\end{figure}

%-----------------------------------------------------------------------
\section{Análisis de Resultados}
%-----------------------------------------------------------------------

\subsection{Calidad en entrenamiento}

\textbf{Greedy y BL empatan en IBEX 35.}
Ambos obtienen el mismo fitness de entrenamiento ($-0{,}2128$) y el mismo
resultado en 2025, lo que indica que el greedy construye directamente el
óptimo local al que converge la BL.
Esto se confirma en el número de evaluaciones: la BL usa solo $1\,852$ de
las $10\,000$ disponibles antes de agotar el vecindario sin encontrar mejora,
lo que significa que parte de una solución aleatoria y en menos de 2\,000
movimientos ya ha llegado a un óptimo local de calidad idéntica a la del greedy.

\textbf{S\&P 100: Greedy y BL idénticos.}
Con $N=97$, ambos algoritmos obtienen prácticamente el mismo fitness de entrenamiento ($0{,}867$).
La BL converge al mismo óptimo local que calcula directamente el greedy, usando las
$10\,000$ evaluaciones disponibles.

\textbf{S\&P 500: Greedy supera a BL en entrenamiento ($1{,}043$ vs $0{,}949$).}
Este resultado, a primera vista sorprendente, se explica por el tamaño del
vecindario: con $N=457$ hay $457 \times 456 = 208\,392$ pares posibles.
En un presupuesto de $10\,000$ evaluaciones, la BL solo puede explorar menos
del $5\,\%$ del vecindario antes de agotarlos.
Partiendo de una solución aleatoria de baja calidad ($\sim 0{,}241$ como
muestra Random), la BL asciende cuanto puede, pero no llega al nivel que
alcanza el greedy usando información global (beneficio acumulado y covarianza
media) para construir su solución en un único paso.
Esto ilustra que para instancias grandes, una heurística bien diseñada puede
superar a la búsqueda local cuando el presupuesto de evaluaciones es pequeño
en relación al tamaño del espacio.

Random es consistentemente el peor en entrenamiento en los tres conjuntos,
lo que era esperable: generar $10\,000$ soluciones aleatorias sin explotar la
estructura del problema es la referencia inferior.

\subsection{Generalización (datos de 2025)}

El comportamiento en test revela un fenómeno notable de \textbf{sobreajuste
progresivo} con el tamaño del problema:

\begin{itemize}
  \item \textbf{IBEX 35}: todos los algoritmos generalizan bien.
    Fitness de 2025 positivo y beneficio positivo para los tres, sin diferencias relevantes.

  \item \textbf{S\&P 100}: Greedy y BL presentan fitness negativo en 2025
    ($-0{,}0407$ y $-0{,}0393$), mientras que \textbf{Random obtiene el único
    fitness positivo} ($+0{,}0065$) y un beneficio bruto casi
    cinco veces mayor ($0{,}0518$ frente a $0{,}0099$--$0{,}0111$).

  \item \textbf{S\&P 500}: los tres algoritmos tienen fitness negativo en 2025,
    y Greedy y BL incluso pierden dinero en bruto (beneficio $-0{,}0633$ y
    $-0{,}048$ respectivamente), mientras que Random logra un beneficio
    marginalmente positivo ($+0{,}006$).
\end{itemize}

La razón por la que Random generaliza mejor que los algoritmos más
sofisticados es contraintuitiva pero estructural: el greedy y la BL
construyen portfolios fuertemente optimizados para la covarianza de 2015--2024.
La alta penalización del riesgo ($\lambda = 500$) amplifica este efecto: los algoritmos
se dedican a minimizar el riesgo histórico, concentrándose en activos con bajo riesgo
relativo en 2015--2024, que en 2025 (un año con condiciones de mercado distintas)
resultaron ser peores. Random, al no explotar esa estructura, genera portfolios más
diversificados que no dependen de patrones históricos específicos y por eso transfieren mejor.

\subsection{Variabilidad entre semillas}

Como es determinista, el \textbf{Greedy produce exactamente el mismo resultado}
en las 50 ejecuciones (misma solución, mismo fitness).
El boxplot del greedy es por tanto un único punto.

La \textbf{BL en IBEX 35} presenta también variabilidad muy baja:
parte de puntos distintos pero converge al mismo óptimo local
(mismo fitness medio que el greedy), con pequeñas variaciones debidas al orden
de exploración.
En S\&P 100 y S\&P 500, la BL no converge antes del presupuesto,
por lo que su resultado final depende en mayor medida del punto de partida;
la varianza es moderada.

\textbf{Random} presenta la mayor varianza en los tres conjuntos,
especialmente pronunciada en S\&P 500 donde el espacio es enorme y
$10\,000$ muestras representan una fracción ínfima.

\subsection{Coste computacional}

\textbf{Dominancia de \texttt{fitness}.}
Cada evaluación de fitness requiere $O(N^2)$ operaciones (producto $w^\top\Sigma w$).
Con $10\,000$ evaluaciones, el coste teórico es $O(10\,000 \cdot N^2)$ operaciones.

\begin{table}[H]
\centering
\footnotesize
\caption{Tiempo medio por semilla (s) con compilación optimizada (-O3, sin OpenMP)}
\label{tab:tiempos}
\begin{tabular}{lrrrrr}
\toprule
 & \textbf{IBEX~35} ($N=30$)
 & \textbf{S\&P~100} ($N=97$)
 & \textbf{S\&P~500} ($N=457$)
 & \textbf{Factor} & \textbf{Teórico} \\
\midrule
Greedy & $0{,}000$ & $0{,}000$ & $0{,}000$ & --- & --- \\
Random & $0{,}020$   & $0{,}050$   & $1{,}193$  & $59{,}7\times$ & $(4{,}7)^2 \approx 22$ \\
BL     & $0{,}001$   & $0{,}029$   & $0{,}745$  & $744{,}3\times$ & $22$ \\
\bottomrule
\end{tabular}
\end{table}

Random y BL escalan de forma casi consistente ($60\times$ y $744\times$ entre IBEX 35 y S\&P 500).
En S\&P 500 ambos hacen las $10\,000$ evaluaciones completas,
donde el tiempo está dominado por \texttt{fitness}: $1{,}193$ s/semilla para Random
y $0{,}745$ s/semilla para BL, una diferencia del $37\%$.

Destaca que \textbf{en IBEX 35 la BL es 20 veces más rápida que Random}
en tiempo promedio ($0{,}001$ s frente a $0{,}020$ s):
la convergencia temprana ($1\,852$ evaluaciones frente a $10\,000$)
reduce el coste real a solo un $18{,}5\,\%$ del presupuesto.

El \textbf{Greedy} es el más eficiente con diferencia: su coste es $O(N^2)$
por el precómputo de heurísticas, pero solo requiere una única
evaluación de fitness, resultando en tiempos despreciables incluso para
$N=457$.

%-----------------------------------------------------------------------

\end{document}
